% 05_syscall
\section{系统调用入口开销}

系统调用入口开销指用户态陷入内核、执行完毕再返回这条路径上，每次调用都要支付的固定成本，与调用具体做什么、是否触及调度或内存无关。一次系统调用无论是读一个字节还是取一个进程号，都要走同一段陷入与返回：用户态执行 \texttt{SVC} 指令触发同步异常陷入 EL1，内核保存现场、按调用号分发到处理函数、处理函数返回后再恢复现场、\texttt{eret} 返回用户态。这段路径的耗时是所有系统调用共同承担的固定成本，调用体本身越轻，它在总耗时中的占比越高。调度在结构与单核算力上追平 Linux 之后，剩余的 IPC 差距集中于此处。

最纯粹的探针是空系统调用 getpid，其处理体只取回一个已在内存中的进程号，几乎不做任何工作，测得的时延近乎全部来自陷入与返回路径。\texttt{syscost.c} 测得 StarryOS 的 getpid 需 \nGetpidBefore{} ns，同机 Linux 仅 \nGetpidLin{} ns，相差约 7 倍。一个不涉及调度、不涉及内存、仅进出内核一趟的调用仍慢 7 倍，说明差距落在这条路径本身。窗口化的 pipe-pong（\texttt{ppong.c}）给出同一结论：两侧均以批处理摊薄单次开销，StarryOS 批处理后的地板仍为每条消息 13.4 µs，Linux 为 1.87 µs。本节先剖析这条路径的构成，指出每次调用都要支付而多数调用并不使用的三处成本，再说明各处的机制与对应优化，最后给出结果。

\subsection{陷入路径的构成}

在 aarch64 上，一次系统调用的时延由两部分构成。第一部分是陷入往返本身，即 \texttt{uctx.run()} 一整趟：\texttt{SVC} 进入、保存通用寄存器现场、按调用号查表分发、处理函数执行、恢复现场、\texttt{eret} 返回。这一段由 axcpu 的架构代码实现，是 ArceOS、StarryOS、Axvisor 三套系统共用的底座；getpid 这类空调用的处理体只占其中极小一部分，其时延几乎全部是这趟往返。第二部分是返回路径上的一组固定记账，在处理函数返回之后、恢复用户态现场之前执行，用于维护跟踪、定时器与 CPU 时间等每进程状态。StarryOS 此前在这段返回路径上串接了三处工作，每次系统调用返回都完整执行一遍，而绝大多数调用并不需要其中任何一处。getpid 的 \nGetpidBefore{} ns 正是这两部分之和：陷入往返的架构地板，叠加三处每调用记账。三处记账构成架构地板之上唯一可摘除的冗余，后续优化集中于此。

\subsection{每次调用固定支付的三处成本}

以下三处均为 StarryOS 此前的既有行为，逐一说明其作用，以及在 getpid 路径上冗余的原因。

\paragraph{ptrace 停点}
ptrace 是进程调试与系统调用跟踪机制，被跟踪进程在每次系统调用的进入与退出处停下，把控制权交给跟踪者（调试器或 strace）读取或修改其状态。StarryOS 此前在每次系统调用返回时都执行一遍这套停点判断与相关状态检查，即便当前进程未被任何跟踪者附加。对未被跟踪的进程，这段逻辑的结论恒为无需停下，其执行不产生任何效果。

\paragraph{POSIX 定时器轮询}
POSIX 定时器（\texttt{timer\_create}、\texttt{setitimer} 等）到期时须向进程投递信号，其到期检查此前由 \texttt{poll\_process\_timer} 在每次系统调用返回时执行。这段检查的代价并不小：它先获取全局 \texttt{PROCESS\_TABLE} 的 SpinRwLock 读锁，再对当前进程做一次 Arc 升级，随后获取该进程 posix\_timers 的 Mutex，最后遍历一遍存放定时器的 BTreeMap。即便进程未装配任何 POSIX 定时器，这一整套加锁、升级、再加锁、遍历仍照常走完，遍历的是一棵空树。对不使用 POSIX 定时器的负载，全程无一处能产生到期事件，冗余在于每次调用都重建并检查一份恒为空的状态。

\paragraph{utime/stime 记账}
CPU 时间记账把每个线程消耗的时间区分为用户态（utime）与内核态（stime）两部分，供 \texttt{getrusage}、\texttt{times} 与 \texttt{/proc} 读取。其实现要求在每次进出内核的边界上完成一次 User 与 Kernel 之间的状态迁移并累加相应时段，而承载状态的 \texttt{thr.time} 为跨 CPU 共享，迁移须在一把 SMP 锁的保护下进行。这把锁架设于系统调用边界之上，每次进出内核各获取一次，锁本身即为这处成本的主体，锁内承载的时间累加相形之下可忽略。

\subsection{三处优化}

针对上述三处成本各实现一处优化，每项上线前均经过一轮对抗评审。三者只移除冗余，不改变对外可见的语义。

\paragraph{ptrace 快路径}
将每次返回的停点判断\ours{收敛到一个预先算好的 \texttt{is\_ptraced} 标志}，对应 Linux 以 \texttt{TIF\_SYSCALL\_WORK} 位门控系统调用跟踪的做法：进程未被跟踪时整段 ptrace 停点逻辑不再进入，被跟踪时才照常执行。getpid 由 \nGetpidBefore{} ns 降至约 900 ns，降幅约 25\%。

\paragraph{poll\_process\_timer 跳过}
在 \texttt{poll\_process\_timer} 之前\ours{引入一个无锁的 \texttt{PosixTimerTable::\allowbreak count} 计数作为门槛}。进程未装配任何 POSIX 定时器时该计数为零，前述的全局读锁、Arc 升级、Mutex 与 BTreeMap 遍历整段跳过，仅在确有定时器时才付出加锁与遍历的代价。此项在隔离测量下将 getpid 由 1026.5 ns 降至 853.3 ns，降幅约 17\%。

\paragraph{无锁 timer\_state 粗粒度 tick 记账}
这是三者中收益最大的一处。先以二分定位成本：将 \texttt{set\_timer\_state} 改为空转，getpid 落到 458 ns，据此测得两次记账调用合计约 395 ns，占当时剩余开销的 46\%。再区分是锁还是算术：仅按时钟 tick 做粗粒度记账、但仍保留那把 SMP 锁的版本只节省约 60 ns，确认成本主要在锁本身，锁内承载的时间累加无足轻重。据此把 User 与 Kernel 两个状态从受锁保护的 \texttt{thr.time} \ours{迁移至每线程一个 \texttt{AtomicU8}}：系统调用边界退化为一次 Relaxed 存储，不获取任何锁、不做记账，完整的时间轮询仅在该线程确有 itimer 装配（\texttt{itimer\_armed}）时才执行。SMP 锁由此彻底移出系统调用边界。粗粒度记账改由时钟 tick 承担，rusage 的 stime 秒占比读到 0.71，与 Linux 的 0.72 对齐，精度对目标负载足够。7 组对抗评审得到 0 个确认缺陷。getpid \keyres{由 853 ns 降至 \nGetpid{} ns，降幅约 50\%}，此时已贴近陷入往返的架构地板。

\subsection{tickacct 精确记账特性}

记账口径本身另经一轮重做，形成 tickacct 特性。它把 utime/stime 的推进由每次系统调用的 \texttt{poll()} 迁移至时钟 tick 与上下文切换，对应 Linux 的 \texttt{TICK\_CPU\_ACCOUNTING}，为此\ours{新增一个全 CPU 的 \texttt{on\_tick} 钩子}。一次 15 组评审查出 5 个确认缺陷，修复其中 2 个 HIGH：一处陈旧的 \texttt{last\_wall\_ns} 会误触发刚装配的 setitimer，一处把系统调用前的用户态计算错记为 stime。评审同时暴露 \texttt{thr.time} 一处既有的跨 CPU 数据竞争，此前以 \texttt{AssumeSync<RefCell<...>>} 掩盖，现以 \texttt{SpinNoIrq<TimeManager>} 收口，并令 \texttt{poll} 返回 \texttt{[Option<Signo>;3]}，使锁不跨信号投递持有；7 组死锁评审得到 0 缺陷，上板确认无死锁。

该特性默认关闭。此前观察到的约 11\% 降幅经核实为发射闭包带来的测量假象，SMP 安全重构之后关闭路径同样取得这份收益，故 tickacct 对 getpid 路径实为中性。已上线的无锁粗粒度 tick 计时给出的秒占比 0.71 已贴合 Linux；tickacct 额外提供的只是全 CPU 一致的精确记账，其每 tick 与上下文切换钩子带来维护与运行成本，在目标负载上并非必要，故保持默认关闭，作为需要精确 rusage 时的可选项保留。

\subsection{结果}

三处优化叠加，\keyres{getpid 由 \nGetpidBefore{} ns 压至 \nGetpid{} ns，相对 Linux 由约 7 倍收敛至 \nGetpidRatio{} 倍}，已并入板级出厂配置。

\figph{../figures/out/F07.png}{getpid 三步累计从 \nGetpidBefore{} ns 降到 \nGetpid{} ns，为 Linux \nGetpidLin{} ns 的 \nGetpidRatio{} 倍，残余逼近 \texttt{uctx.run} 的陷入下限}

\begin{table}[H]\centering\small
\begin{tabular}{@{}lrr@{}}
\toprule
版本 & getpid (ns) & 相对 Linux \\
\midrule
Linux & \nGetpidLin{} & 1× \\
StarryOS 基线 & \nGetpidBefore{} & 7.2× \\
加 ptrace 快路径 & 900 & 5.4× \\
加 poll\_process\_timer 跳过 & 853 & 5.1× \\
二分参考（set\_timer\_state 空转） & 458 & 2.7× \\
无锁 timer\_state 粗粒度 tick（上线） & \nGetpid{} & \nGetpidRatio{}× \\
精确记账回退（非 tickacct） & 975 & 5.8× \\
\bottomrule
\end{tabular}
\caption{getpid 微基准逐层下降，板级 p50（中间两级为实现当期的单次测量），二分参考行为定位记账成本的中间探测点（\texttt{set\_timer\_state} 空转），精确记账回退行给出需要精确记账时的可选口径，高于上线值，不属该下降序列}
\end{table}

残余的 \nGetpid{} ns 已贴近前述陷入往返的架构地板，包含 \texttt{SVC} 进入、保存现场、按调用号分发、恢复现场与 \texttt{eret} 返回。该路径属架构层，由 ArceOS、StarryOS、Axvisor 三套系统共用，其中不存在可摘除的冗余，改动波及面大且无对应收益，故不触碰。pipe 数据路径是下一步的优化方向，\texttt{pipe\_wr} 约 5 µs 对应 Linux 的约 0.4 µs。
